\chapter{Introduzione}
\label{cap:introduzione}

\section{Contesto e motivazioni}
\label{sec:contesto-e-motivazioni}

La simulazione è uno degli strumenti utilizzati per studiare sistemi formati
da molte entità che interagiscono fra loro e con lo spazio in cui si trovano:
reti di dispositivi, flotte di veicoli, gruppi di persone che si spostano in
un'area urbana. Quando lo scenario di riferimento è collocato in un luogo reale,
l'attendibilità di un esperimento non dipende solamente dal comportamento
dei suoi attori, ma anche da quanto fedelmente è riprodotto l'ambiente in
cui essi agiscono.

Di quell'ambiente fanno parte grandezze che variano da un punto all'altro e
lungo il tempo, e che non vengono prodotte dagli attori: la portata di un fiume,
la temperatura dell'aria, la concentrazione di inquinanti in atmosfera, il pericolo
d'incendio di una zona boschiva. Chi si trova sul territorio non può controllarle,
le percepisce, e vi reagisce prendendo decisioni: cambiando percorso, diramando un
allarme, coordinandosi con altri attori. Queste misure costituiscono
lo stimolo esterno rispetto al quale il comportamento degli attori di una
simulazione prende significato.

Grandezze di questo genere, per buona parte del pianeta, sono misurate,
ricostruite e pubblicate con continuità. Copernicus\footnote{Il programma Copernicus: \url{https://www.copernicus.eu/it}.},
il programma di osservazione della Terra dell'Unione Europea, distribuisce liberamente archivi
che coprono decenni di storia, i quali spaziano dalle rianalisi climatiche ai
sistemi di allerta alluvionale. Eventi realmente accaduti dispongono quindi
di descrizioni quantitative pubbliche, consultabili rispetto a istanti temporali
e posizioni diverse.

Che quei dati siano pubblici e accessibili non basta però perché un attore,
nel corso di una simulazione, possa chiedere quanto valga una grandezza
nel punto in cui si trova. I dati sono pubblicati in formati binari specialistici,
organizzati secondo convenzioni della comunità scientifica che li produce,
e vanno richiesti a servizi remoti che li preparano su richiesta, impiegando
un tempo variabile e senza garanzia di riuscita. Soprattutto, sono campionati,
quindi noti solamente in corrispondenza dei punti di una griglia e di una
successione di istanti temporali, mentre la posizione di un attore, e il
momento in cui esso interroga l'ambiente, sono quantità continue che un
simulatore misura a partire dall'inizio della simulazione.

Alchemist~\cite{alchemist}, il simulatore di sistemi complessi sviluppato
presso l'Università di Bologna su cui si aggiunge questo lavoro, possiede già
la maggior parte di ciò che serve per collocare un esperimento in un luogo preciso:
sa caricare cartografia reale, disporvi gli attori e vincolarne il movimento
alla rete stradale. Ciò che non possiede è un modo per far percepire loro una
misurazione campionata in una posizione geografica. Le grandezze diffuse
nell'ambiente sono rappresentate da un'astrazione dedicata, il \emph{layer}, le cui
realizzazioni disponibili producono però valori calcolati da regole fissate
al momento della costruzione. Una simulazione può quindi essere accurata nel
collocare i propri nodi ma non altrettanto in ciò che essi percepiscono una
volta collocati. 

\section{Il contributo di questa tesi}
\label{sec:contributo}

Il lavoro qui descritto colma questa distanza con un nuovo modulo del
simulatore, che rende i dati distribuiti da Copernicus interrogabili dagli
attori come grandezza diffusa nell'ambiente, nel punto in cui l'attore
si trova e nell'istante di simulazione in cui lo chiede.

Per fare in modo che ciò sia possibile occorre leggere file il cui schema non
è noto a priori e ricondurli a una rappresentazione uniforme; produrre un valore
anche dove una misura non è disponibile, secondo un criterio che non può essere
unico perché dipende da che cosa la grandezza rappresenta; e ottenere i dati
dai servizi Copernicus senza alcun intervento manuale, conservandoli in locale
in modo che esecuzioni successive non debbano richiederli una seconda volta.
Il tutto senza modificare le astrazioni del modello del simulatore, e in modo
dichiarabile direttamente nel file di configurazione di una simulazione,
senza scrivere codice.

Il codice sorgente del modulo, nella versione documentata in questa tesi, è
archiviato in modo permanente e liberamente
consultabile\footnote{Archivio permanente del modulo: \url{https://archive.softwareheritage.org/swh:1:dir:1246f419d7f827ab6a7767d13f922af4845632c9;origin=https://github.com/Wanes01/Alchemist;visit=swh:1:snp:1145cad987574aa5003aa00de817ba00591905e3;anchor=swh:1:rev:58f72d98e0b10428b192d8e555a357f0da7366c6;path=/alchemist-geospatial/}}.

\section{Struttura della tesi}
\label{sec:struttura-della-tesi}

La tesi è strutturata nei seguenti capitoli:

\begin{itemize}
  \item \textbf{Contesto} (\Cref{chap:contesto}): vengono descritti gli elementi su cui il
  lavoro poggia: le parti del simulatore inerenti al progetto, i servizi attraverso
  cui Copernicus distribuisce i propri dati, i formati e le convenzioni con cui
  questi sono codificati e lo standard a cui si ispirano le richieste remote.

  \item \textbf{Analisi} (\Cref{chap:analisi}): si scompone il problema nelle sue parti e
  se ne ricavano i requisiti che la soluzione deve soddisfare, vengono discusse le questioni che
  vanno delimitate prima di poter essere tradotte in scelte progettuali e si individuano le
  entità in gioco.

  \item \textbf{Progettazione} (\Cref{chap:progettazione}): viene esposta la struttura che
  deriva dall'analisi, ovvero la suddivisione del modulo, le astrazioni che compongono
  ciascuna delle sue parti e le decisioni che le attraversano.

  \item \textbf{Implementazione} (\Cref{chap:implementazione}): ci si sofferma sui punti
  in cui la realizzazione ha richiesto scelte non deducibili dalla progettazione, imposte
  dal formato dei dati, dal comportamento dei servizi remoti o dal meccanismo con
  cui Alchemist istanzia i componenti di una simulazione.

  \item \textbf{Validazione} (\Cref{chap:validazione}): vengono riportate le evidenze
  del corretto funzionamento del modulo, le quali comprendono la suite di test
  automatici, le verifiche condotte al di fuori di essi e una simulazione dimostrativa
  che impiega l'intera catena su dati reali, costruita attorno agli eventi alluvionali
  che hanno colpito la Romagna nel maggio 2023.

  \item \textbf{Conclusioni} (\Cref{chap:conclusioni}): viene tracciato un bilancio
  del lavoro rispetto agli obiettivi, vengono dichiarati i limiti del lavoro e
  si indicano le direzioni di sviluppo che la struttura realizzata rende possibili. 
\end{itemize}
